\section*{Ethics Statement}

We mine public repository histories, and the only personal data we touch is
authorship metadata. Each commit in a file's history carries an author name and
email; we reduce them at extraction to a count of distinct editors, the
covariate in \autoref{sec:hazard}. No artifact we release carries an identity.
We redistribute no repository content, only derived tables and the code that
rebuilds them (\apxref{sec:appendix:artifacts}). These files are public, but
nobody wrote them for this measurement. Every result we report is therefore an
aggregate over \FieldRepos{} repositories rather than a judgment about any one
maintainer.

Our own recommendation carries the main risk. We tell maintainers to delete
instructions whose rationale they can recover, and an operator who automates
that rule will delete instructions whose rationale was real. Our protocol
empties a nontrivial share of prompts, and the uncommented arm scores higher on
exactly those worlds (\apxref{sec:appendix:calibration}). Writing a comment
removes nothing and is therefore safe. Acting on one is not. An operator
deploying the protocol should keep a person in the deletion path and hold
safety-relevant instructions out of scope until someone tests the protocol on
them.

We did not audit the corpus for offensive content, and we print only three
verbatim strings, all comments from our own run
(\autoref{fig:comment_examples}). We recruited no participants, and an author
produced the one hand-annotation here (\autoref{tab:gates}); an independent
re-annotation would fix that. We train nothing, and \apxref{sec:appendix:compute} reports
the compute.
